\section{Conclusion}
\label{sec:conclusion}

We introduced \dsName, a comprehensive benchmark for evaluating GPS reasoning in LLMs, spanning geometric coordinate computation and applied geographic reasoning tasks. Our evaluation reveals that GPS reasoning remains partially solved: models perform reasonably on basic geometric operations but struggle with complex spherical geometry and fine-grained place association. We find that world knowledge does not transfer to coordinate computation, with most models showing stronger applied reasoning than pure GPS skills. Geographic bias is task-specific---applied tasks exhibit substantial regional disparities driven by sparse coordinate-to-place mappings, while pure computation remains region-agnostic. Models encode coarse geography but lack fine-grained localization, and robustness to coordinate noise confirms genuine understanding rather than memorization. Finetuning reveals capability trade-offs where geometric computation improves at the expense of world knowledge.

Our findings suggest that applications requiring precise geospatial reasoning should not assume LLMs can reliably perform geodetic calculations, and developers should evaluate regional performance before deploying location-aware systems. The success of GPS augmentation on downstream benchmarks indicates practical value in providing explicit coordinates for geographic queries. Future work should explore training methods that improve geodetic computation without catastrophic forgetting, extend evaluation to multimodal settings with maps and satellite imagery, and investigate retrieval-augmented approaches for fine-grained place association. \dsName and all evaluation code are publicly available to support reproducible research.
